%% file attn.tex
%% author thegreatchaos@thegreatchaos.com
%% date 2026/06/11 16:09:33
\subsubsection{代码结构分析}
\begin{enumerate}
    \item build/CMakeLists.txt. 忽略, 把代码嵌入进去或者让ai来做
    \item FMHA
\begin{lstlisting}
cd src
grep -irl fmha 得到以下
flash_attn/flash_api.cpp  // mha_varlen_fwd
xpu/attn/xe_2/chunk_prefill_configure.cmake //中定义了fmha_forward_configure函数, 会被xpu/attn/xe_2/CMakeLists.txt:5调用来依据chunk_prefill_kernel生成不同配置的.cpp
xpu/attn/xe_2/chunk_prefill_utils.hpp //用来dispatch不同的配置的Kernel
xpu/attn/xe_2/paged_decode_xe2.cpp //cutlass_paged_decode_xe2的实现
xpu/attn/xe_2/fmha_xe2.cpp //cutlass_chunk_prefill_xe2的实现
xpu/attn/xe_2/kernel/chunk_prefill_kernel.hpp //FMHAProblemShape, XeFMHAFwdKern两class/struct的定义
xpu/attn/xe_2/kernel/paged_decode_kernel.hpp //DecodeProblemShape, XeFMHAFwdSplitKVKernel, ReduceSplitK class/struct的定义
xpu/attn/xe_2/collective/chunk_prefill_epilogue.hpp //FMHAFwdEpilogue, DecodeFwdEpilogue class定义
xpu/attn/xe_2/collective/chunk_prefill_mainloop.hpp //FMHAFwdMainLoop, DecodeFwdMainLoop class定义
xpu/attn/xe_2/collective/chunk_prefill_scheduler.hpp //XeFHMAIndividualTileScheduler, DecodeTileScheduler, XeReduceSplitKTileScheduler class定义
xpu/attn/xe_2/paged_decode.hpp//多种decode policy, {DecodeKernelLauncher, PageDecodeConfig}结构体定义, decode_policy_dispatch_impl的实现
xpu/attn/xe_2/chunk_prefill.hpp //KernelLauncher, FMHAConfig结构体定义, policy_dispatch_impl的实现
xpu/attn/xe_2/paged_decode_utils.hpp //utils: decode_policy_dispatch_func, dispatch_by_{head_size, page_size}, cutlass_paged_decode_impl
xpu/attn/xe_2/CMakeLists.txt
xpu/attn/attn_interface.cpp //cutlass_{chunk_prefill, paged_decode}_interface的实现
\end{lstlisting}
    \item 调用链路
\begin{lstlisting}
flash_attn_varlen_func(){ //vllm_xpu_kernels/flash_attn_interface.py
    ...
    out, softmax_lse = torch.ops._vllm_fa2_C.varlen_fwd(){ //vllm_xpu_kernels/flash_attn_interface.py
	//从此开始C/C++
	mha_varlen_fwd(){ //csrc/flash_attn/flash_api.cpp, 文件尾部注册算子mha_varlen_fwd
	    //q_type check,目前只支持BF16&Half 
	    bool is_paged = block_table_.has_value();o
	    bool is_varlen = true;//后续全部按照varlen处理
	    if (max_seqlen_q > 1 || !is_paged) {//prefill分支
		cutlass_chunk_prefill_interface(){ //csrc/xpu/attn/attn_interface.cpp
		    cutlass_chunk_prefill_xe2(){ //csrc/xpu/attn/xe_2/fmha_xe2.cpp
			cutlass_chunk_prefill_impl(){ //同上文件
			    //按照head_size dispatch 到不同的policy 模板特化 chunk_policy_head{64,96,128,192,256,512}[_b16]
			    //模板定义在 csrc/xpu/attn/xe_2/fmha_utils.hpp中struct chunk_policy_head{64..512}_b16
			    //实例化在 csrc/xpu/attn/xe_2/chunk_prefill_extern.hpp. 将不同配置的policy实例化成extern template vod policy_dispatch_impl<...>(sycl::queue&, CutlassQKType& const chunk_prefill_args_t&);
			}
		    }
		}
	    }else{//decode分支,  paged + 单token
		get_num_splits(); //计算kv split数
		cutlass_paged_decode_interface(){//csrc/xpu/attn/attn_interface.cpp
		    cutlass_paged_decode_xe2{// csrc/xpu/attn/xe_2/paged_decode_xe2.cpp
			cutlass_paged_decode_impl(){ //同上文件
			    dispatch_by_page_size<QGroup>(){ //QGroup=8,16, csrc/xpu/attn/xe_2/paged_decode_utils.hpp
				dispatch_by_head_size<QGroup, PageSize>(){//文件同上, PageSize=16,32,64
				    decode_plociy_dispatch_func<>(){//文件同上
					decode_policy_dispatch_impl<>(){//csrc/xpu/attn/xe_2/paged_decode.hpp
					    return PageDecodeConfig<不同配置policy>::kernel_dispatch(){ //文件同上
						PageDecodeConfig::run(){//文件同上
						    DecodeKernelLauncher<FMHAKernel, ReduceSplitKernel, VarLen> launcher;
						    launcher.run(queue, args, hw_info){
							//cutlass
							compat::experimental::laucnh<...>()
						    }
						}
					    }
					}
				    }
				}
			    }
			}
		    }
		}
	    }
	}
    }
    ...
}
\end{lstlisting}
\end{enumerate}
